\begin{description}
	\item[1行目] 先ほどの推定に使ったデータファイルです。
	\item[2行目] \verb|lag|関数で，データを1行ずらした列を作ります。
	\item[3行目] 体重を記録した日の変数\verb|date|と，先ほど作った一行ずらした変数\verb|lag|はいずれも日付に関する変数ですので，日付型に変更します。
	\item[4行目] 日付変数の引き算をして，一行下のデータと何日ずれていたのか算出しています。
	\item[5行目] 1日以上ずれている日をフィルタリングで抜き出しています。
\end{description}
